% Source: https://artofproblemsolving.com/community/c3h1454046_the_four_cyclists_a_curious_problem
% Stable identity: aops-topic-1454046-q01
A rider at $v$ miles per hour completes a lap of $\dfrac{1}{3}$ mile in $\dfrac{1}{3v}$ hours, and is at the centre exactly at the multiples of that time. The lap times are
\[
\dfrac{1}{18},\ \dfrac{1}{27},\ \dfrac{1}{36},\ \dfrac{1}{45}\text{ hours}=200,\ 133\dfrac{1}{3},\ 100,\ 80\text{ seconds}.
\]
In units of $\dfrac{1}{3}$ second these are $600,400,300,240$. All four are at the centre at a time $T$ exactly when $T$ is a common multiple of these, and the least common multiple is $\operatorname{lcm}(600,400,300,240)=1200$ units, i.e.\ $400$ seconds $=6$ minutes $40$ seconds. (In that time the riders complete $2$, $3$, $4$ and $5$ laps.) So they are all together at the centre every $6$ min $40$ s, and the fourth time after the start is after $4\cdot400=1600$ seconds $=26$ min $40$ s, at
\[
\boxed{12{:}26{:}40}.
\]
(If the start is counted as the first meeting, the fourth meeting is the third after the start, at $12{:}20$.)
